\section{Discussion}

In this section, we discuss the broader implications of \textbf{Dynamic Activity-Dependent Pruning (DADP)}, analyze how its reverse Hebbian mechanism bridges neuroscience and machine learning, and address potential limitations alongside future research directions.

\subsection{Neuroscience Connections \& Reverse Hebbian Plasticity}

Classical biological learning is dominated by Hebbian plasticity (``cells that fire together, wire together''). However, mammalian neurodevelopment relies heavily on activity-dependent synapse elimination during childhood and adolescence \citep{faust2021mechanisms}. Synapses that remain idle or transmit weak error signals are gradually pruned to streamline brain connectivity. 

DADP operationalizes this biological mechanism in deep neural networks via the activation-derivative product $I_{ij}^l = \mathbb{E}\left[\left| a_i^{l-1} \cdot \frac{\partial \mathcal{L}}{\partial y_j^l} \right|\right]$. Our empirical findings demonstrate that reverse Hebbian pruning naturally resolves long-standing challenges in artificial network compression:
\begin{enumerate}
    \item \textbf{Organic Structural Allocation}: Unlike rigid dynamic sparse training algorithms (e.g., RigL) that enforce uniform or pre-specified layer budgets, DADP allows the network to discover its own optimal capacity distribution organically.
    \item \textbf{Emergent Structural Protection}: DADP automatically identifies critical architectural pathways—such as downsample skip connections in ResNet-18—and preserves them without manual intervention or specialized regularizers.
\end{enumerate}

\subsection{Limitations \& Future Directions}

While DADP achieves strong performance across diverse architectures and vision/NLP tasks, two primary limitations present promising avenues for future work:

\begin{enumerate}
    \item \textbf{Training Set Dependency \& Generalization Boundary}: Because DADP importance scores are calculated on training sample loss gradients $\frac{\partial \mathcal{L}_{\text{train}}}{\partial y_j^l}$, extreme pruning threshold settings ($\tau \ge \tau_{\text{crit}}$) risk overfitting to training set feature patterns. Incorporating validation set sensitivity estimates or dropout-based stochastic gradient sampling could improve generalization near extreme capacity boundaries.
    \item \textbf{Hardware Acceleration for Unstructured Masks}: During early training epochs, DADP operates as a soft binary mask ($W \odot M$). While emergent filter deletion allows physical tensor compression at step boundaries, maximum memory and speed gains on specialized hardware (e.g., NVIDIA Tensor Cores) will require custom 2:4 sparse kernel integrations.
\end{enumerate}
